%% fig_dial_example.tex
%% Dial のアルゴリズム：リンク尤度・前進処理・後退処理の例示
%%
%% ネットワーク：r(Origin), A, B, s(Destination)
%% リンクコスト：t_{rA}=1, t_{rB}=3, t_{As}=2, t_{AB}=1, t_{Bs}=1
%% 最短経路費用（Dijkstra）：c(r)=0, c(A)=1, c(B)=2, c(s)=3
%% θ=1 でのリンク尤度 L[i→j] = exp(θ{c(j)−c(i)−t_{ij}})：
%%   L[r→A]=1.00, L[r→B]=e^{-1}≈0.37, L[A→s]=1.00, L[A→B]=1.00, L[B→s]=1.00
%%
%% Usage: %% fig_dial_example.tex
%% Dial のアルゴリズム：リンク尤度・前進処理・後退処理の例示
%%
%% ネットワーク：r(Origin), A, B, s(Destination)
%% リンクコスト：t_{rA}=1, t_{rB}=3, t_{As}=2, t_{AB}=1, t_{Bs}=1
%% 最短経路費用（Dijkstra）：c(r)=0, c(A)=1, c(B)=2, c(s)=3
%% θ=1 でのリンク尤度 L[i→j] = exp(θ{c(j)−c(i)−t_{ij}})：
%%   L[r→A]=1.00, L[r→B]=e^{-1}≈0.37, L[A→s]=1.00, L[A→B]=1.00, L[B→s]=1.00
%%
%% Usage: %% fig_dial_example.tex
%% Dial のアルゴリズム：リンク尤度・前進処理・後退処理の例示
%%
%% ネットワーク：r(Origin), A, B, s(Destination)
%% リンクコスト：t_{rA}=1, t_{rB}=3, t_{As}=2, t_{AB}=1, t_{Bs}=1
%% 最短経路費用（Dijkstra）：c(r)=0, c(A)=1, c(B)=2, c(s)=3
%% θ=1 でのリンク尤度 L[i→j] = exp(θ{c(j)−c(i)−t_{ij}})：
%%   L[r→A]=1.00, L[r→B]=e^{-1}≈0.37, L[A→s]=1.00, L[A→B]=1.00, L[B→s]=1.00
%%
%% Usage: %% fig_dial_example.tex
%% Dial のアルゴリズム：リンク尤度・前進処理・後退処理の例示
%%
%% ネットワーク：r(Origin), A, B, s(Destination)
%% リンクコスト：t_{rA}=1, t_{rB}=3, t_{As}=2, t_{AB}=1, t_{Bs}=1
%% 最短経路費用（Dijkstra）：c(r)=0, c(A)=1, c(B)=2, c(s)=3
%% θ=1 でのリンク尤度 L[i→j] = exp(θ{c(j)−c(i)−t_{ij}})：
%%   L[r→A]=1.00, L[r→B]=e^{-1}≈0.37, L[A→s]=1.00, L[A→B]=1.00, L[B→s]=1.00
%%
%% Usage: \input{assets/assignment/fig_dial_example.tex}

\begin{tikzpicture}[
    every node/.style={font=\small},
    nd/.style={circle, draw, thick, minimum size=0.75cm,
               font=\small\bfseries, inner sep=1pt, fill=white},
    cnd/.style={rectangle, draw, thick, rounded corners=3pt,
                minimum width=0.80cm, minimum height=0.80cm,
                font=\small\bfseries, inner sep=4pt, fill=gray!10},
    arr/.style={->, >=stealth, thick},
    arrfaint/.style={->, >=stealth, thick, gray!50},
    lbl/.style={font=\footnotesize, fill=white, inner sep=1.5pt},
    lbll/.style={font=\scriptsize, fill=white, inner sep=1pt, text=blue!70!black},
    lbllt/.style={font=\scriptsize, fill=white, inner sep=1pt, text=gray!60}
]

%% ===== ノード（最短費用 c(i) を表示） =====
\node[cnd, label={[font=\scriptsize,gray]below:$c=0$}] (r) at (0.0,  0.0) {$r$};
\node[nd,  label={[font=\scriptsize,gray]below=12pt, right=-10pt:$c=1$}] (A) at (3.0,  1.5) {$A$};
\node[nd,  label={[font=\scriptsize,gray]below:$c=2$}] (B) at (3.0, -1.5) {$B$};
\node[cnd, label={[font=\scriptsize,gray]below:$c=3$}] (s) at (6.0,  0.0) {$s$};

%% ===== リンク（コスト+尤度を2段表記） =====
%% r→A（L=1.00，最短経路上）
\draw[arr] (r)
    to node[lbl, above, sloped] {$t=1$}
    node[lbll, below, sloped] {$L=1.00$}
    (A);

%% r→B（L=e^{-1}≈0.37，非最短）
\draw[arrfaint] (r)
    to node[lbl, below, sloped] {$t=3$}
    node[lbllt, above, sloped] {$L=e^{-1}{\approx}0.37$}
    (B);

%% A→B（L=1.00，最短経路上）
\draw[arr] (A)
    to node[lbl, right] {$t=1$}
    node[lbll, left] {$L=1.00$}
    (B);

%% A→s（L=1.00，最短経路上）
\draw[arr] (A)
    to node[lbl, above, sloped] {$t=2$}
    node[lbll, below, sloped] {$L=1.00$}
    (s);

%% B→s（L=1.00，最短経路上）
\draw[arr] (B)
    to node[lbl, below, sloped] {$t=1$}
    node[lbll, above, sloped] {$L=1.00$}
    (s);

%% ===== OD需要 =====
\draw[<->, >=stealth, thick, gray, dashed, out=70, in=110]
    (r) to node[above=3pt, font=\footnotesize, gray] {$q_{rs}=10$} (s);

%% ===== 前進/後退パス方向の矢印（説明ラベル） =====
\node[draw, rounded corners=2pt, font=\footnotesize, align=left,
      inner sep=5pt, fill=white]
  at (3.0, -3.25) {
    {\bfseries 前進処理（Step 3）}：$c(i)$ の昇順（$r \to s$ 方向）にノードを処理\\
    {\bfseries 後退処理（Step 4）}：$c(i)$ の降順（$s \to r$ 方向）にノードを処理\\[1pt]
    \textcolor{blue!70!black}{青字}：$L>0$（合理的リンク）\quad
    \textcolor{gray!50}{灰色}：$L<1$（非最短方向）
  };

\end{tikzpicture}


\begin{tikzpicture}[
    every node/.style={font=\small},
    nd/.style={circle, draw, thick, minimum size=0.75cm,
               font=\small\bfseries, inner sep=1pt, fill=white},
    cnd/.style={rectangle, draw, thick, rounded corners=3pt,
                minimum width=0.80cm, minimum height=0.80cm,
                font=\small\bfseries, inner sep=4pt, fill=gray!10},
    arr/.style={->, >=stealth, thick},
    arrfaint/.style={->, >=stealth, thick, gray!50},
    lbl/.style={font=\footnotesize, fill=white, inner sep=1.5pt},
    lbll/.style={font=\scriptsize, fill=white, inner sep=1pt, text=blue!70!black},
    lbllt/.style={font=\scriptsize, fill=white, inner sep=1pt, text=gray!60}
]

%% ===== ノード（最短費用 c(i) を表示） =====
\node[cnd, label={[font=\scriptsize,gray]below:$c=0$}] (r) at (0.0,  0.0) {$r$};
\node[nd,  label={[font=\scriptsize,gray]below=12pt, right=-10pt:$c=1$}] (A) at (3.0,  1.5) {$A$};
\node[nd,  label={[font=\scriptsize,gray]below:$c=2$}] (B) at (3.0, -1.5) {$B$};
\node[cnd, label={[font=\scriptsize,gray]below:$c=3$}] (s) at (6.0,  0.0) {$s$};

%% ===== リンク（コスト+尤度を2段表記） =====
%% r→A（L=1.00，最短経路上）
\draw[arr] (r)
    to node[lbl, above, sloped] {$t=1$}
    node[lbll, below, sloped] {$L=1.00$}
    (A);

%% r→B（L=e^{-1}≈0.37，非最短）
\draw[arrfaint] (r)
    to node[lbl, below, sloped] {$t=3$}
    node[lbllt, above, sloped] {$L=e^{-1}{\approx}0.37$}
    (B);

%% A→B（L=1.00，最短経路上）
\draw[arr] (A)
    to node[lbl, right] {$t=1$}
    node[lbll, left] {$L=1.00$}
    (B);

%% A→s（L=1.00，最短経路上）
\draw[arr] (A)
    to node[lbl, above, sloped] {$t=2$}
    node[lbll, below, sloped] {$L=1.00$}
    (s);

%% B→s（L=1.00，最短経路上）
\draw[arr] (B)
    to node[lbl, below, sloped] {$t=1$}
    node[lbll, above, sloped] {$L=1.00$}
    (s);

%% ===== OD需要 =====
\draw[<->, >=stealth, thick, gray, dashed, out=70, in=110]
    (r) to node[above=3pt, font=\footnotesize, gray] {$q_{rs}=10$} (s);

%% ===== 前進/後退パス方向の矢印（説明ラベル） =====
\node[draw, rounded corners=2pt, font=\footnotesize, align=left,
      inner sep=5pt, fill=white]
  at (3.0, -3.25) {
    {\bfseries 前進処理（Step 3）}：$c(i)$ の昇順（$r \to s$ 方向）にノードを処理\\
    {\bfseries 後退処理（Step 4）}：$c(i)$ の降順（$s \to r$ 方向）にノードを処理\\[1pt]
    \textcolor{blue!70!black}{青字}：$L>0$（合理的リンク）\quad
    \textcolor{gray!50}{灰色}：$L<1$（非最短方向）
  };

\end{tikzpicture}


\begin{tikzpicture}[
    every node/.style={font=\small},
    nd/.style={circle, draw, thick, minimum size=0.75cm,
               font=\small\bfseries, inner sep=1pt, fill=white},
    cnd/.style={rectangle, draw, thick, rounded corners=3pt,
                minimum width=0.80cm, minimum height=0.80cm,
                font=\small\bfseries, inner sep=4pt, fill=gray!10},
    arr/.style={->, >=stealth, thick},
    arrfaint/.style={->, >=stealth, thick, gray!50},
    lbl/.style={font=\footnotesize, fill=white, inner sep=1.5pt},
    lbll/.style={font=\scriptsize, fill=white, inner sep=1pt, text=blue!70!black},
    lbllt/.style={font=\scriptsize, fill=white, inner sep=1pt, text=gray!60}
]

%% ===== ノード（最短費用 c(i) を表示） =====
\node[cnd, label={[font=\scriptsize,gray]below:$c=0$}] (r) at (0.0,  0.0) {$r$};
\node[nd,  label={[font=\scriptsize,gray]below=12pt, right=-10pt:$c=1$}] (A) at (3.0,  1.5) {$A$};
\node[nd,  label={[font=\scriptsize,gray]below:$c=2$}] (B) at (3.0, -1.5) {$B$};
\node[cnd, label={[font=\scriptsize,gray]below:$c=3$}] (s) at (6.0,  0.0) {$s$};

%% ===== リンク（コスト+尤度を2段表記） =====
%% r→A（L=1.00，最短経路上）
\draw[arr] (r)
    to node[lbl, above, sloped] {$t=1$}
    node[lbll, below, sloped] {$L=1.00$}
    (A);

%% r→B（L=e^{-1}≈0.37，非最短）
\draw[arrfaint] (r)
    to node[lbl, below, sloped] {$t=3$}
    node[lbllt, above, sloped] {$L=e^{-1}{\approx}0.37$}
    (B);

%% A→B（L=1.00，最短経路上）
\draw[arr] (A)
    to node[lbl, right] {$t=1$}
    node[lbll, left] {$L=1.00$}
    (B);

%% A→s（L=1.00，最短経路上）
\draw[arr] (A)
    to node[lbl, above, sloped] {$t=2$}
    node[lbll, below, sloped] {$L=1.00$}
    (s);

%% B→s（L=1.00，最短経路上）
\draw[arr] (B)
    to node[lbl, below, sloped] {$t=1$}
    node[lbll, above, sloped] {$L=1.00$}
    (s);

%% ===== OD需要 =====
\draw[<->, >=stealth, thick, gray, dashed, out=70, in=110]
    (r) to node[above=3pt, font=\footnotesize, gray] {$q_{rs}=10$} (s);

%% ===== 前進/後退パス方向の矢印（説明ラベル） =====
\node[draw, rounded corners=2pt, font=\footnotesize, align=left,
      inner sep=5pt, fill=white]
  at (3.0, -3.25) {
    {\bfseries 前進処理（Step 3）}：$c(i)$ の昇順（$r \to s$ 方向）にノードを処理\\
    {\bfseries 後退処理（Step 4）}：$c(i)$ の降順（$s \to r$ 方向）にノードを処理\\[1pt]
    \textcolor{blue!70!black}{青字}：$L>0$（合理的リンク）\quad
    \textcolor{gray!50}{灰色}：$L<1$（非最短方向）
  };

\end{tikzpicture}


\begin{tikzpicture}[
    every node/.style={font=\small},
    nd/.style={circle, draw, thick, minimum size=0.75cm,
               font=\small\bfseries, inner sep=1pt, fill=white},
    cnd/.style={rectangle, draw, thick, rounded corners=3pt,
                minimum width=0.80cm, minimum height=0.80cm,
                font=\small\bfseries, inner sep=4pt, fill=gray!10},
    arr/.style={->, >=stealth, thick},
    arrfaint/.style={->, >=stealth, thick, gray!50},
    lbl/.style={font=\footnotesize, fill=white, inner sep=1.5pt},
    lbll/.style={font=\scriptsize, fill=white, inner sep=1pt, text=blue!70!black},
    lbllt/.style={font=\scriptsize, fill=white, inner sep=1pt, text=gray!60}
]

%% ===== ノード（最短費用 c(i) を表示） =====
\node[cnd, label={[font=\scriptsize,gray]below:$c=0$}] (r) at (0.0,  0.0) {$r$};
\node[nd,  label={[font=\scriptsize,gray]below=12pt, right=-10pt:$c=1$}] (A) at (3.0,  1.5) {$A$};
\node[nd,  label={[font=\scriptsize,gray]below:$c=2$}] (B) at (3.0, -1.5) {$B$};
\node[cnd, label={[font=\scriptsize,gray]below:$c=3$}] (s) at (6.0,  0.0) {$s$};

%% ===== リンク（コスト+尤度を2段表記） =====
%% r→A（L=1.00，最短経路上）
\draw[arr] (r)
    to node[lbl, above, sloped] {$t=1$}
    node[lbll, below, sloped] {$L=1.00$}
    (A);

%% r→B（L=e^{-1}≈0.37，非最短）
\draw[arrfaint] (r)
    to node[lbl, below, sloped] {$t=3$}
    node[lbllt, above, sloped] {$L=e^{-1}{\approx}0.37$}
    (B);

%% A→B（L=1.00，最短経路上）
\draw[arr] (A)
    to node[lbl, right] {$t=1$}
    node[lbll, left] {$L=1.00$}
    (B);

%% A→s（L=1.00，最短経路上）
\draw[arr] (A)
    to node[lbl, above, sloped] {$t=2$}
    node[lbll, below, sloped] {$L=1.00$}
    (s);

%% B→s（L=1.00，最短経路上）
\draw[arr] (B)
    to node[lbl, below, sloped] {$t=1$}
    node[lbll, above, sloped] {$L=1.00$}
    (s);

%% ===== OD需要 =====
\draw[<->, >=stealth, thick, gray, dashed, out=70, in=110]
    (r) to node[above=3pt, font=\footnotesize, gray] {$q_{rs}=10$} (s);

%% ===== 前進/後退パス方向の矢印（説明ラベル） =====
\node[draw, rounded corners=2pt, font=\footnotesize, align=left,
      inner sep=5pt, fill=white]
  at (3.0, -3.25) {
    {\bfseries 前進処理（Step 3）}：$c(i)$ の昇順（$r \to s$ 方向）にノードを処理\\
    {\bfseries 後退処理（Step 4）}：$c(i)$ の降順（$s \to r$ 方向）にノードを処理\\[1pt]
    \textcolor{blue!70!black}{青字}：$L>0$（合理的リンク）\quad
    \textcolor{gray!50}{灰色}：$L<1$（非最短方向）
  };

\end{tikzpicture}
